% GENERATED FILE. Edit canonical lessons, then run npm run generate:books.
% canonical-source-hash: fnv1a64:3618852b7c8a262f
% canonical-lessons: TE-C234-kuppa, TE-C234-eruvu, TE-C234-peda, TE-C234-nagali, TE-C234-kodavali, TE-R234-first-pass-words-from-chapters-226-229, TE-R234-second-pass-words-from-chapters-230-234

\chapter{Heap, Fertilizer, Cow dung, Plough, Sickle}
\label{ch:te-heap-fertilizer-cow-dung-plough-sickle}
\hlchaptermodality{234}

\begin{chapteropening}
\textbf{By the end of this chapter:} \emph{I can say a heap, fertilizer, cow dung, a plough and a sickle.}

Complete the last lesson of chapter 234: I can say a heap, fertilizer, cow dung, a plough and a sickle.
\end{chapteropening}

\section[\texorpdfstring{\te{కుప్ప} (kuppa)}{కుప్ప (kuppa)}]{\te{కుప్ప} (kuppa) — a heap, a pile}

\label{lesson:TE-C234-kuppa}



\begin{quote}
Before the new one: say the Telugu for soot, then the Telugu for charcoal.
\end{quote}

\subsection*{You'll want to know: \te{కుప్ప}}
\textbf{\te{కుప్ప}} — \emph{kuppa} — \textquotedblleft{}a heap, a pile\textquotedblright{}.

\begin{grammarlens}[title={What you've built}]
One more word for this chapter.
\end{grammarlens}

\subsection*{Your turn}
\begin{itemize}
\raggedright
  \item \emph{Say it:} \emph{kuppa}
  \item \emph{Say it:} \emph{kuppa}, once more
  \item \emph{Say it:} say \emph{kuppa}
  \item \emph{Recall:} say the Telugu for soot, then the Telugu for charcoal, then say \emph{kuppa} again
\end{itemize}

\begin{tcolorbox}[breakable,colback=teal!4,colframe=teal!35!black,title={Before you move on}]
What does \te{కుప్ప} mean? (\textquotedblleft{}A heap, a pile\textquotedblright{}.) Say it once more.
\end{tcolorbox}
\section[\texorpdfstring{\te{ఎరువు} (eruvu)}{ఎరువు (eruvu)}]{\te{ఎరువు} (eruvu) — fertilizer, manure}

\label{lesson:TE-C234-eruvu}



\begin{quote}
Before the new one: say the Telugu for charcoal, then the Telugu for a heap.
\end{quote}

\subsection*{You'll want to know: \te{ఎరువు}}
\textbf{\te{ఎరువు}} — \emph{eruvu} — \textquotedblleft{}fertilizer, manure\textquotedblright{}.

\begin{grammarlens}[title={What you've built}]
One more word for this chapter.
\end{grammarlens}

\subsection*{Your turn}
\begin{itemize}
\raggedright
  \item \emph{Say it:} \emph{eruvu}
  \item \emph{Say it:} \emph{eruvu}, once more
  \item \emph{Say it:} say \emph{eruvu}
  \item \emph{Recall:} say the Telugu for charcoal, then the Telugu for a heap, then say \emph{eruvu} again
\end{itemize}

\begin{tcolorbox}[breakable,colback=teal!4,colframe=teal!35!black,title={Before you move on}]
What does \te{ఎరువు} mean? (\textquotedblleft{}Fertilizer, manure\textquotedblright{}.) Say it once more.
\end{tcolorbox}
\section[\texorpdfstring{\te{పేడ} (pēḍa)}{పేడ (pēḍa)}]{\te{పేడ} (pēḍa) — cow dung}

\label{lesson:TE-C234-peda}



\begin{quote}
Before the new one: say the Telugu for a heap, then the Telugu for fertilizer.
\end{quote}

\subsection*{You'll want to know: \te{పేడ}}
\textbf{\te{పేడ}} — \emph{pēḍa} — \textquotedblleft{}cow dung\textquotedblright{}.

\begin{grammarlens}[title={What you've built}]
One more word for this chapter.
\end{grammarlens}

\subsection*{Your turn}
\begin{itemize}
\raggedright
  \item \emph{Say it:} \emph{pēḍa}
  \item \emph{Say it:} \emph{pēḍa}, once more
  \item \emph{Say it:} say \emph{pēḍa}
  \item \emph{Recall:} say the Telugu for a heap, then the Telugu for fertilizer, then say \emph{pēḍa} again
\end{itemize}

\begin{tcolorbox}[breakable,colback=teal!4,colframe=teal!35!black,title={Before you move on}]
What does \te{పేడ} mean? (\textquotedblleft{}Cow dung\textquotedblright{}.) Say it once more.
\end{tcolorbox}
\section[\texorpdfstring{\te{నాగలి} (nāgali)}{నాగలి (nāgali)}]{\te{నాగలి} (nāgali) — a plough}

\label{lesson:TE-C234-nagali}



\begin{quote}
Before the new one: say the Telugu for fertilizer, then the Telugu for cow dung.
\end{quote}

\subsection*{You'll want to know: \te{నాగలి}}
\textbf{\te{నాగలి}} — \emph{nāgali} — \textquotedblleft{}a plough\textquotedblright{}.

\begin{grammarlens}[title={What you've built}]
One more word for this chapter.
\end{grammarlens}

\subsection*{Your turn}
\begin{itemize}
\raggedright
  \item \emph{Say it:} \emph{nāgali}
  \item \emph{Say it:} \emph{nāgali}, once more
  \item \emph{Say it:} say \emph{nāgali}
  \item \emph{Recall:} say the Telugu for fertilizer, then the Telugu for cow dung, then say \emph{nāgali} again
\end{itemize}

\begin{tcolorbox}[breakable,colback=teal!4,colframe=teal!35!black,title={Before you move on}]
What does \te{నాగలి} mean? (\textquotedblleft{}A plough\textquotedblright{}.) Say it once more.
\end{tcolorbox}
\section[\texorpdfstring{\te{కొడవలి} (koḍavali)}{కొడవలి (koḍavali)}]{\te{కొడవలి} (koḍavali) — a sickle}

\label{lesson:TE-C234-kodavali}



\begin{quote}
Before the new one: say the Telugu for cow dung, then the Telugu for a plough.
\end{quote}

\subsection*{You'll want to know: \te{కొడవలి}}
\textbf{\te{కొడవలి}} — \emph{koḍavali} — \textquotedblleft{}a sickle\textquotedblright{}.

\begin{grammarlens}[title={What you've built}]
One more word for this chapter.
\end{grammarlens}

\subsection*{Your turn}
\begin{itemize}
\raggedright
  \item \emph{Say it:} \emph{koḍavali}
  \item \emph{Say it:} \emph{koḍavali}, once more
  \item \emph{Say it:} say \emph{koḍavali}
  \item \emph{Recall:} say the Telugu for cow dung, then the Telugu for a plough, then say \emph{koḍavali} again
\end{itemize}

\begin{tcolorbox}[breakable,colback=teal!4,colframe=teal!35!black,title={Before you move on}]
What does \te{కొడవలి} mean? (\textquotedblleft{}A sickle\textquotedblright{}.) Say it once more.
\end{tcolorbox}
\section[(dialogue)]{Words from chapters 226-229 — a first pass}

\label{lesson:TE-R234-first-pass-words-from-chapters-226-229}



\begin{quote}
Say the words for a plough and a sickle.
\end{quote}

\subsection*{Your turn}
Say these aloud:

\begin{itemize}
\raggedright
  \item a map, a document, a file, a form, a card
  \item a young man, elders, neighbours, an enemy, a bridegroom
  \item a crab, a scorpion, a bat, a nest, a tail
  \item firecrackers, a decorated performing bull seen at Sankranti, sunshine, a raindrop, fog
\end{itemize}

\begin{tcolorbox}[breakable,colback=teal!4,colframe=teal!35!black,title={Before you move on}]
A map? (\textbf{\te{పటం}}, \emph{paṭaṁ}.) A document? (\textbf{\te{పత్రం}}, \emph{patraṁ}.) A file? (\textbf{\te{ఫైలు}}, \emph{phailu}.) A form? (\textbf{\te{ఫారం}}, \emph{phāraṁ}.) A card? (\textbf{\te{కార్డు}}, \emph{kārḍu}.) A young man? (\textbf{\te{యువకుడు}}, \emph{yuvakuḍu}.) Elders? (\textbf{\te{పెద్దలు}}, \emph{peddalu}.) Neighbours? (\textbf{\te{ఇరుగుపొరుగు}}, \emph{iruguporugu}.) An enemy? (\textbf{\te{శత్రువు}}, \emph{śatruvu}.) A bridegroom? (\textbf{\te{పెళ్ళికొడుకు}}, \emph{peḷḷikoḍuku}.) A crab? (\textbf{\te{పీత}}, \emph{pīta}.) A scorpion? (\textbf{\te{తేలు}}, \emph{tēlu}.) A bat? (\textbf{\te{గబ్బిలం}}, \emph{gabbilaṁ}.) A nest? (\textbf{\te{గూడు}}, \emph{gūḍu}.) A tail? (\textbf{\te{తోక}}, \emph{tōka}.) Firecrackers? (\textbf{\te{టపాకాయలు}}, \emph{ṭapākāyalu}.) A decorated performing bull seen at Sankranti? (\textbf{\te{గంగిరెద్దు}}, \emph{gaṅgireddu}.) Sunshine? (\textbf{\te{ఎండ}}, \emph{eṇḍa}.) A raindrop? (\textbf{\te{చినుకు}}, \emph{cinuku}.) Fog? (\textbf{\te{పొగమంచు}}, \emph{pogamañcu}.)
\end{tcolorbox}
\section[(dialogue)]{Words from chapters 230-234 — a second pass}

\label{lesson:TE-R234-second-pass-words-from-chapters-230-234}



\begin{quote}
Say the words for a poem and a painting.
\end{quote}

\subsection*{Your turn}
Say these aloud:

\begin{itemize}
\raggedright
  \item a poem, a painting, a camera, entertainment, an actor
  \item a fan, pliers, a ladder, glue, a ruler
  \item a jar, a sack, a tub, a barrel, a mug
  \item a shoulder cloth, a bubble, rust, soot, charcoal
  \item a heap, fertilizer, cow dung, a plough, a sickle
\end{itemize}

\begin{tcolorbox}[breakable,colback=teal!4,colframe=teal!35!black,title={Before you move on}]
A poem? (\textbf{\te{కవిత}}, \emph{kavita}.) A painting? (\textbf{\te{చిత్రం}}, \emph{citraṁ}.) A camera? (\textbf{\te{కెమెరా}}, \emph{kemerā}.) Entertainment? (\textbf{\te{వినోదం}}, \emph{vinōdaṁ}.) An actor? (\textbf{\te{నటుడు}}, \emph{naṭuḍu}.) A fan? (\textbf{\te{ఫ్యాను}}, \emph{phyānu}.) Pliers? (\textbf{\te{పటకారు}}, \emph{paṭakāru}.) A ladder? (\textbf{\te{నిచ్చెన}}, \emph{niccena}.) Glue? (\textbf{\te{జిగురు}}, \emph{jiguru}.) A ruler? (\textbf{\te{కొలబద్ద}}, \emph{kolabadda}.) A jar? (\textbf{\te{జాడీ}}, \emph{jāḍī}.) A sack? (\textbf{\te{బస్తా}}, \emph{bastā}.) A tub? (\textbf{\te{తొట్టి}}, \emph{toṭṭi}.) A barrel? (\textbf{\te{పీపా}}, \emph{pīpā}.) A mug? (\textbf{\te{మగ్గు}}, \emph{maggu}.) A shoulder cloth? (\textbf{\te{కండువా}}, \emph{kaṇḍuvā}.) A bubble? (\textbf{\te{బుడగ}}, \emph{buḍaga}.) Rust? (\textbf{\te{తుప్పు}}, \emph{tuppu}.) Soot? (\textbf{\te{మసి}}, \emph{masi}.) Charcoal? (\textbf{\te{బొగ్గు}}, \emph{boggu}.) A heap? (\textbf{\te{కుప్ప}}, \emph{kuppa}.) Fertilizer? (\textbf{\te{ఎరువు}}, \emph{eruvu}.) Cow dung? (\textbf{\te{పేడ}}, \emph{pēḍa}.) A plough? (\textbf{\te{నాగలి}}, \emph{nāgali}.) A sickle? (\textbf{\te{కొడవలి}}, \emph{koḍavali}.)
\end{tcolorbox}
